\documentclass[11pt]{article}
% Shared preamble for all solution files.
\usepackage[margin=1.1in]{geometry}
\usepackage{enumitem}
\usepackage{tikz}
\usepackage[hidelinks]{hyperref}
\usepackage{amsmath,amssymb,amsthm}
\newcommand{\fa}{\mathfrak a}
\newcommand{\fb}{\mathfrak b}
\newcommand{\fc}{\mathfrak c}
\newcommand{\fp}{\mathfrak p}
\newcommand{\fq}{\mathfrak q}
\newcommand{\fm}{\mathfrak m}
\newcommand{\fn}{\mathfrak n}
\newcommand{\fr}{\mathfrak r}
\newcommand{\Nil}{\mathfrak N}
\newcommand{\Jac}{\mathfrak R}
\newcommand{\rad}{r}
\newcommand{\Spec}{\operatorname{Spec}}
\newcommand{\Max}{\operatorname{Max}}
\newcommand{\Ann}{\operatorname{Ann}}
\newcommand{\Ker}{\operatorname{Ker}}
\newcommand{\Z}{\mathbf Z}
\newcommand{\Q}{\mathbf Q}
\newcommand{\R}{\mathbf R}
\newcommand{\C}{\mathbf C}
\newcommand{\ci}{\mathrm{c}}
\newcommand{\ex}{\mathrm{e}}

\setlength{\parindent}{0pt}
\setlength{\parskip}{0.55em}
% \exercise{label}{statement}  -- followed by the solution text
\newcommand{\exercise}[2]{\par\bigskip\noindent{\large\textbf{Exercise #1}}\par\smallskip
  \begin{quote}\small\textbf{Problem.}\ #2\end{quote}\par}
\newcommand{\pf}{\noindent\textbf{Solution.}\ }

\usepackage{needspace}
\AddToHook{cmd/section/before}{\Needspace{14\baselineskip}}
\let\chapterexercise\exercise
\renewcommand{\exercise}[2]{\Needspace{8\baselineskip}\chapterexercise{#1}{#2}}
\title{Liu, Chapter 10: Reduction of algebraic curves\\[2pt]\large Solutions and corrections to the exercises}
\author{}\date{}
\begin{document}\maketitle
\textbf{Conventions.} The 57 labels follow the four exercise lists on printed pp.~479--483, 503--505, 528--532 and 552--556. Write $R$ for a DVR, $t$ for a uniformizer, $K$ for its fraction field and $k$ for its residue field. Models are proper unless an open or formal model is explicitly under discussion. Semistability itself does not imply properness in Definition 10.3.14. All assertions about geometric components or their dual graphs are made after a separable splitting extension or over the specified algebraically closed residue field. A loop contributes twice to a graph's valence. Degrees include residue-field degrees.

\textbf{Methods and corrections.} We use the chapter's reduction criteria, relative duality, minimal-model uniqueness and stable-reduction theorem, stating explicitly where the latter is used. Solution 2.5(c) uses the established Neron mapping property; 3.16 explicitly uses the Jacobian specialization theorem and Deligne--Mumford's Jacobian criterion. The original formulations of 1.6, 1.10, 1.12(b), 1.13, 1.15(b), 1.17(a),(d), 3.3, 3.5--7, 3.8(b1), 3.13 and 4.6 need correction or additional hypotheses. Examples, not merely revised claims, are given below. In 4.4(c) the presentation arrow is reversed in print. In 4.5 the generic constant field matters. Exercise 4.7 retains the algebraically closed residue field and connected-fibre conventions of Subsection 10.4.2; the changed block over a nonrational center is also computed. No assertion of global good reduction is inferred just from a local coordinate calculation.

\textbf{Additional reference for 3.16.} The proof there deliberately uses the Jacobian form of the stable-reduction criterion, rather than claiming it follows just by omitting multiplicities in Proposition 3.44. The input is Deligne--Mumford, \emph{The irreducibility of the space of curves of given genus}, Publications Math. IHES 36 (1969), Theorem 2.4, together with the Picard/Neron specialization theorem. The other solutions use the indicated results from Liu and explicitly displayed local calculations.

\section*{10.1 Models and reductions}

\exercise{1.1}{Extend Proposition 1.8 from affine Dedekind bases to normal algebraic curves over a field.}
\pf Such a base $S$ is regular and quasi-projective. Spread a projective embedding and equations of the smooth generic curve over a dense affine open $U\subset S$, and shrink $U$ so the resulting family is smooth and projective. At each of the finitely many missing points, choose a projective local model as in Proposition 1.8. Its equations and the generic identification spread to a neighborhood of that point. Shrink the overlaps so these families agree with the chosen smooth model over $U$. Gluing gives a proper flat model over $S$; normalization and the surface resolution theorem give a regular one. It is projective: a proper regular surface over a quasi-projective curve is projective over that curve, by the surface projectivity theorem used in 8.3.16, or by extending an ample divisor from the generic curve and adding vertical divisors. Now contract exceptional curves for a relatively minimal model. For genus at least one, Chapter 9 gives the unique minimal regular model; resolving fibre crossings gives the normal-crossings versions. This proves all assertions of Proposition 1.8 in the new setting. \hfill$\square$

\exercise{1.2}{Construct a model over affine $S$ for a geometrically reduced normal curve over $K(S)$.}
\pf As in the definition of model, the curve must be integral and projective; without properness of the generic curve it cannot be the generic fibre of a projective model. Take a flat projective closure of its equations over $S$ and normalize it. The finiteness theorem for normalization with geometrically reduced generic fibre (Theorem 8.2.39) applies, even if $S$ was not assumed excellent. The normalization is therefore finite over the closure and projective over $S$. It is integral and dominant, hence flat over the Dedekind base, and its generic fibre is the original normal curve. If ``curve'' is allowed nonproper here, the same construction models its projective completion; the original curve has only an open model in that projective surface. \hfill$\square$

\exercise{1.3}{Show that good reduction persists under change of Dedekind base.}
\pf A dominant morphism $S'\to S$ is required to define the displayed extension of function fields. Base change a smooth proper model near a given point $s$ to a neighborhood of any $s'$ above it. Smoothness, properness and flatness survive base change, and the generic fibre is $C_{K(S')}$. Thus it is a smooth proper model at $s'$. If the convention requires integrality, use the connected component belonging to the chosen generic curve; for geometrically connected $C$ the whole base change is integral. \hfill$\square$

\exercise{1.4}{For the Klein quartic $x^3y+y^3z+z^3x=0$, prove normality of its integral model and good reduction away from seven.}
\pf Put $F=x^3y+y^3z+z^3x$. A coordinate-zero point of the quartic is one of the three coordinate vertices, where some partial derivative is one. At a singular geometric point with all coordinates nonzero set $A=x^3y$, $B=y^3z$, $C=z^3x$. Multiplication of the three derivative equations by $x,y,z$ gives
\[ 3A+C=0,\qquad A+3B=0,\qquad B+3C=0. \]
For characteristic three these force $A=B=C=0$, impossible. Otherwise they imply $28A=0$; adding them and using $A+B+C=0$ rules out characteristic two as well (in characteristic two the equations give $A=B=C$, and their sum on the quartic is $A\ne0$). Thus only characteristic seven can be singular. In that characteristic the ratios of the coordinates solving these equations form a finite set, so the special fibre is generically smooth and in particular reduced. The integral hypersurface is Cohen--Macaulay; its possible nonregular points form a finite subset of that fibre, of codimension two in the surface. Serre's criterion makes it normal. The same model is smooth and proper at every prime other than seven, proving good reduction there. \hfill$\square$

\exercise{1.5}{Over a complete DVR with separably closed residue field, compare rational points and multiplicity-one geometrically reduced components in some model.}
\pf If a model has such a component, choose a nonempty open of it where the component and the morphism are smooth and no other component occurs. A smooth variety over a separably closed field has a rational point in every nonempty open: take an etale chart and a separable closed point. Smooth Hensel lifting gives an $R$-section through that point and hence a $K$-point.

Conversely let $P\in C(K)$. Resolve a model strongly; its generic fibre is the given normal curve, and the section defined by properness meets the regular surface. At its closed point, the retraction to the local DVR of the section shows $t\notin\fm^2$. Thus the fibre is regular there with multiplicity one and rational residue field. A regular curve over a field is smooth at a rational regular point. The component through it is generically smooth and therefore geometrically reduced. This proves the other implication. Equivalently one may first spread the smooth rational point of the regular generic curve and then resolve away from it. \hfill$\square$

\exercise{1.6}{Test the asserted equivalence between good reduction of a smooth conic and splitting over an etale DVR extension.}
\pf \textbf{The converse is false in genus zero.} Take $R=k[[t]]$, $\operatorname{char}k\ne2$, and a nonsquare $a\in k^*$. The conic
\[ u^2-av^2+tw^2=0 \]
has a point over the finite etale extension $k(\sqrt a)[[t]]$. Its displayed model is regular: at the vertex of the special fibre the equation has linear term $t$. That fibre is integral over $k$ and singular, so the model is relatively minimal by Exercise 9.3.1. It has no smooth model over $R$: the elementary-transformation classification in that exercise permits a change between relatively minimal conic models only at a split smooth fibre. Equivalently the associated quaternion class $(a,t)$ has nonzero residue and cannot extend to an Azumaya conic over $R$.

Good reduction does imply etale-local splitting: choose a separable point on the smooth special conic and lift it through an etale neighborhood by the smooth local-section lemma. The correct equivalence with etale-local splitting is \emph{semistable} reduction, proved in 3.8. The counterexample above already has a geometric nodal special fibre and hence is semistable. \hfill$\square$

\exercise{1.7}{Give the genus-one version of Lemma 1.24 and determine the special fibre of the minimal regular model.}
\pf Let $\Gamma'$ be the smooth genus-at-least-one normalization supplied by that lemma's hypothesis. Its strict transform is not exceptional and survives on the minimal model $X$. Write its image as $\Gamma$. The normalization sequence gives $p_a(\Gamma)\ge p_a(\Gamma')\ge1$. Minimality gives $K_X\cdot\Gamma\ge0$, whereas $K_X\cdot X_s=0$ forces this intersection, and those with all other components, to be zero. Adjunction yields $\Gamma^2=2p_a(\Gamma)-2\ge0$. Negative semidefiniteness forces equality; connectedness of the fibre then forces it to have no other components. Equality in the genus comparison shows that $\Gamma=\Gamma'$ is smooth of genus one. Thus $X_s=d\Gamma$. If $C(K)\ne\varnothing$, its section intersects the fibre in degree one, forcing $d=1$. Without a point multiple fibres can occur, so that last conclusion must not be reversed without further assumptions. \hfill$\square$

\exercise{1.8}{Bound degrees of points in a formal fibre of a regular model.}
\pf For $x\in C^+(\widetilde x)$, apply the horizontal intersection formula to its closure $H$. Its degree is the sum of local intersections with $X_s$, including residue degrees. At $\widetilde x$,
\[ i_{\widetilde x}(H,X_s)\ge\mu_{\widetilde x}(X_s)
=\sum_i d_i\mu_{\widetilde x}(\Gamma_i). \]
The residue degree factor is at least one, proving the stated bound. The stronger bound includes $[k(\widetilde x):k]$, as in Exercise 9.2.10. \hfill$\square$

\exercise{1.9}{Prove the integral hyperelliptic equation criterion for good reduction, including residue characteristic two and the degree-reduction induction.}
\pf (a) A smooth special affine hyperelliptic curve of the full generic genus supplies a multiplicity-one component with smooth projective normalization of that genus. Lemma 1.24 gives good reduction for genus at least two. For genus one use 1.7: the multiplicity-one affine component forces $d=1$. Thus the conclusion also holds there. The smooth hyperelliptic special curve is understood geometrically connected, as in the example.

(b) On the smooth model's inverse image of $\mathbf A^1_R$, use the free rank-two algebra $B=R[v]\oplus R[v]z$ from the example. A generic generator $z_0$ has an equation $z_0^2+G_0z_0=F_0$ with degrees at most $g+1$ and $2g+2$. Scale $z_0$ by a scalar in $K^*$ so that $F_0,G_0$ are integral. Normality of $B$ then gives $z_0\in B$. Since $z$ and $z_0$ are both bases modulo $K[v]$, write $z_0=az+H(v)$ with $0\ne a\in R$ and $H\in R[v]$. If $a$ is a unit we are done. Otherwise reduction gives
\[ \overline H^{,2}+\overline G_0\overline H=\overline F_0. \]
If $\deg\overline H>g+1$, its square has uniquely highest degree, a contradiction. Thus $H=H_0+tH_1$ with $\deg H_0\le g+1$ and $H_1\in R[v]$.

(c) Put $z_1=(z_0-H_0)/t=(a/t)z+H_1\in B$. Its quadratic equation has coefficients
\[ G_1=(G_0+2H_0)/t,\qquad
F_1=(F_0-G_0H_0-H_0^2)/t^2. \]
They lie in $R[v]$: they are the trace and minus norm of the integral element $z_1$, and $R[v]$ is normal. They retain the same degree bounds. The valuation of its leading scalar $a/t$ is one smaller. Iterate to reach a unit scalar. This proves the desired bounded integral equation without completing the square, so also in characteristic two. \hfill$\square$

\exercise{1.10}{Test the equal-root-distance criterion for potential good reduction of a hyperelliptic equation.}
\pf Equal distances of all finite roots are sufficient. After a finite extension, translate by one root and divide $x$ by an element with their common distance valuation. The normalized roots have distinct reductions. Adjoin a square root of the remaining scalar coefficient, so that the equation is integral with separable special polynomial of degree $2g+1$ or $2g+2$. Example 1.26 gives good reduction.

\textbf{Necessity for an arbitrary displayed equation is false: it is not invariant under fractional linear changes of $x$.} Over $\Q_7$ the genus-two curve $y^2=\prod_{a=0}^5(x-a)$ has good reduction. Put $x=8+1/z$ and $Y=yz^3$. Then
\[ Y^2=\prod_{a=0}^5\bigl(1+(8-a)z\bigr). \]
Its roots are $1/(a-8)$. The root for $a=1$ has valuation $-1$, all the others have valuation zero; differences involving that root have valuation $-1$, whereas differences between two of the others have valuation zero. Thus the printed criterion fails even for a curve with good reduction.

The invariant correction is: after a finite extension and a fractional linear coordinate change on $\mathbf P^1$, the \emph{whole branch divisor}, including infinity when appropriate, has pairwise distinct reductions. A smooth model has a smooth quotient by the tame hyperelliptic involution and an etale branch divisor, proving necessity. Conversely that branch configuration gives a smooth double cover after a possible quadratic scalar twist. This is exactly the potential version of Example 1.26. \hfill$\square$

\exercise{1.11}{Describe the minimal reduction of $y^2=x^6+atx^4+bt^2x^2+t^3$, and its good quadratic base change.}
\pf Let $s^2=t$, $x=sX$, $y=s^3Y$. The equation becomes
\[ Y^2=X^6+aX^4+bX^2+1. \]
The discriminant hypothesis and $2\in R^*$ make its special polynomial separable, so its smooth projective model over $R[s]$ has genus two and good reduction. The deck involution acts as $(s,X,Y)\mapsto(-s,-X,-Y)$. On the special curve its only fixed points are the two points at infinity, both rational: in coordinates $q=1/X$, $V=Y/X^3$, they are $q=0,V=\pm1$. The quotient curve has genus one by tame Hurwitz, since the degree-two map from a genus-two curve has exactly two ramification points.

The quotient surface is regular away from these fixed points. At either fixed point the completed action is $(s,q)\mapsto(-s,-q)$; its invariants are $t=s^2$, $z=sq$, $w=q^2$ with $z^2=tw$. This ordinary quadratic surface singularity resolves by one blow-up. The resulting fibre is
\[ X_s=2\Gamma+E_1+E_2, \]
where $\Gamma$ is the smooth elliptic quotient and the $E_i$ are disjoint rational curves meeting it transversally once. The fibre equations give $E_i^2=-2$ and $\Gamma^2=-1$. None is an exceptional divisor: $\Gamma$ has genus one and the rational curves have square $-2$. Thus this is the minimal regular model. Its nonreduced fibre proves bad reduction over $R$, while the displayed substitution proves good reduction over $K(\sqrt t)$. \hfill$\square$

\exercise{1.12}{Spread a smooth stratification in characteristic zero and lift special rational points over completed fraction fields.}
\pf (a) The reduced generic fibre has a finite stratification by smooth locally closed $K$-schemes, since $K$ is perfect. Spread the strata and their immersions to some open $V\subset S$, and shrink $V$ so they remain smooth and cover all points of $X_V$. The omitted constructible set has image missing the generic point, so another shrinking removes it. In the printed union $f_i(V_i)$, the argument should be $X_i$.

(b) For $s\in V$, every $k(s)$-rational point belongs to a unique stratum with the same residue field. Its smooth morphism lifts that point to $\widehat{\mathcal O}_{S,s}$, giving the required $K_s$-rational point of $X$.

\textbf{The extra demand that $V$ contain a closed point is not automatic.} Let $S=\Spec k[[t]]$, $\operatorname{char}k=0$, and $X=\Spec R[z]/(z^2-t)$. This is flat and finite, with a rational point in the special fibre and no $K=k((t))$-point. Any open containing the closed point is all of $S$. Thus the asserted (b) is false for this Dedekind base. If $S$ has infinitely many closed points, the dense open obtained in (a) contains some of them and the intended conclusion holds. \hfill$\square$

\exercise{1.13}{Examine finiteness of closures of generic points over an arbitrary regular base and the claimed reduction map.}
\pf \textbf{Both assertions fail even with the stated flatness hypothesis.} Let $T=\mathbf A^2_k=\Spec k[u,v]$, $X=\mathbf P^1_T$. In the generic fibre take the rational point $[u:v]$. Its closure is the graph closure of this rational map, namely $\operatorname{Bl}_{(u,v)}T$. Its fibre above the origin is $\mathbf P^1$, so the closure is neither finite over $T$ nor isomorphic to it. Different tangent directions approaching the origin give different limiting points, and there is no reduction map of the asserted kind. Here $X\to T$ is smooth, flat and projective.

For a normal \emph{one-dimensional} base the intended assertions are valid. The closure of a closed generic point is an integral proper curve with zero-dimensional fibres, hence finite. If the point is rational, this finite map is birational and normality of the base makes it an isomorphism. These are the DVR and Dedekind cases used earlier in the chapter. For (c), without flatness the blow-up $X=\operatorname{Bl}_0T\to T$ itself, with its unique generic point over $k(T)$, is another counterexample; it is the same graph phenomenon. Flatness cannot rescue the higher-dimensional version. \hfill$\square$

\exercise{1.14}{Define formal fibres for arbitrary finite-type schemes by affine neighborhoods, prove independence, and extend Proposition 1.40.}
\pf If $X$ is affine, take a projective compactification $X\subset Y$. A point of $Y^+(\widetilde x)$, with residue field a finite extension $L/K$, extends to $\Spec R_L\to Y$ with closed point $\widetilde x$. The inverse image of the open $X$ contains the closed point of this local scheme, hence the entire scheme. Thus its generic point lies in $X$. Such extensions in $X$ are unique by separatedness, so the resulting set is independent of the compactification.

If $U,V$ are affine neighborhoods of $\widetilde x$, the same argument applied to $U\cap V$ shows their formal fibres coincide. Therefore the definition for general $X$ is independent of its affine neighborhood, even without assuming $X$ itself separated. More intrinsically, for finite $L/K$ it is the set of continuous local $R$-homomorphisms
\[ \widehat{\mathcal O}_{X,\widetilde x}\longrightarrow R_L. \]
The universal property of completion proves part (a) of Proposition 1.40. If $\widetilde x$ is rational, generators of its special maximal ideal give a presentation by $R[[T_1,\ldots,T_q]]$ and the homomorphisms are exactly the tuples in $\fm_L^q$ satisfying its relations. Smoothness makes this completed ring a power-series ring when $q$ is the relative dimension. An affine closed embedding gives the coordinate version in part (c). These proofs are local and need no projectivity beyond the temporary compactification. \hfill$\square$

\exercise{1.15}{Describe morphisms of formal fibres and the formal group law of a smooth group scheme.}
\pf (a) Pullback gives a continuous local map of completed local rings at $\widetilde y=f(\widetilde x)$ and $\widetilde x$. Composing homomorphisms into $R_L$ defines the required map, compatibly with finite extensions $L/K$.

(b) \textbf{The displayed power-series coordinates require rational special points.} For rational $\widetilde x,\widetilde y$, smoothness identifies their completed rings with $R[[T_1,\ldots,T_d]]$ and $R[[U_1,\ldots,U_e]]$. Write $f^*U_j=F_j(T)$. Their constant terms belong to $tR$, so they converge and map $\fm_L^d$ into $\fm_L^e$. They are defined by a single completed local homomorphism over $R$, hence do not depend on $L$. For a nonrational point the claimed parametrization can fail even in dimension zero: take $X=\Spec R'$ for a nontrivial finite etale extension of a complete DVR. Its special point has no $K$-rational lift, whereas $\fm_K^0$ is a singleton. The corrected version uses its unramified coefficient extension and choices of residue-field embeddings.

(c) The identity is rational, so no correction is needed. Complete the multiplication at $(o,o)$ and obtain $F(T,V)\in R[[T,V]]^d$, with
\[ F(T,0)=T,\quad F(0,V)=V,\quad
F(F(T,V),W)=F(T,F(V,W)). \]
The inverse morphism similarly gives a power-series inverse. These identities follow from the group-scheme identities on completed local rings. Evaluation in any finite extension's maximal ideal gives exactly the kernel-of-reduction group law. The directed union over finite subextensions of a fixed algebraic closure defines the stated formal group of points. For a commutative $G$, $F(T,V)=F(V,T)$. \hfill$\square$

\exercise{1.16}{Describe regular models of $\mathbf P^1$ and the fibre graph when exactly one exceptional divisor remains.}
\pf (a) A relatively minimal conic model has integral fibres. A rational generic point specializes to a smooth rational special point, so each fibre is $\mathbf P^1_k$ by Exercise 9.3.1. Hence the model is smooth. Every regular model dominates such a model by point blow-ups. Blowing up a point of a normal-crossings divisor preserves normal crossings, so all these models have normal crossings in the scheme-theoretic sense.

(b) Reverse a factorization to a smooth model. The last-created exceptional curve $E$ meets one component if its center was a smooth fibre point, and two if its center was an intersection. It has no further modifications because it is the unique exceptional divisor available for the first contraction. The total-transform formula gives its multiplicity as the sum of those one or two component multiplicities. Repeating after contracting it shows that at each stage the preceding centers lie on the exceptional configuration: a blow-up elsewhere would leave two disjoint final exceptional curves. This observation gives the whole graph by reversing the contractions.

(c) More explicitly, the permitted operations on the weighted graph are insertion of a vertex of weight $a$ at a smooth point of a weight-$a$ component, and subdivision of an edge of weights $a,b$ by a vertex of weight $a+b$. Starting with weight one, the single-exceptional condition forces each subsequent operation to destroy the preceding exceptional curve or curves, so all nontrivial branching lies on the path from the final exceptional vertex toward the original weight-one ends. Reversing these operations gives the staircase-shaped tree in the text: the terminal arms have weight one, the weights are nonincreasing along the horizontal chains and strictly decreasing down a branch away from $E$. Equality occurs only at a smooth-point insertion; at an edge subdivision the parent weight $a+b$ strictly exceeds each child. This proves the depicted ordering and the existence of the two weight-one ends, with the original component and the first exceptional component supplying them. \hfill$\square$

\exercise{1.17}{Study normal models of $\mathbf P^1$ over a complete DVR with algebraically closed residue field, and detect regularity by degrees in formal fibres.}
\pf (a) A regular resolution is a tree of projective lines and has $H^1$ of its reduced fibre zero. Contracting connected subtrees gives the reduced fibre of any normal model. The normalization sequence, or Theorem 7.5.19, shows that every component is $\mathbf P^1$ and its singularities are ordinary \emph{multiple} points. The word ``double'' in the printed assertion is too restrictive. For example blow up three distinct points on the fibre of $\mathbf P^1_R$ and contract the strict original fibre, of square $-3$. The three exceptional lines then meet at one ordinary triple point on the normal model. Its completed reduced local ring is the three-axis ring, not a node.

For (b), contract all original components except $\Gamma_1$ and call the resulting one-component model $Y$. This contraction is an isomorphism near $\widetilde x$, since $\widetilde x$ belongs to none of the contracted curves. Let $Z\to Y$ be the minimal resolution. No curve contracted by it can be a Castelnuovo exceptional divisor. Every regular model of $\mathbf P^1$ that is not smooth has such a divisor, so the only possible one on $Z$ is the retained strict transform of $\Gamma_1$; it is unique. Apply 1.16 to it. Its weight $d_1$ is at least two: its canonical intersection is $-1$, all other components have nonnegative canonical intersection, and $K_Z\cdot Z_s=-2$ implies $d_1\ge2$.

If this unique exceptional curve has two neighbors, their weights add to $d_1$ by 1.16(b), and both are smaller. Each adjacent contracted cluster therefore contains a component of weight $<d_1$. If it has one neighbor, the complementary tree is connected and contains the weight-one ends of 1.16(c), again giving a smaller weight. Every nontrivial fibre of $Z\to Y$ is one of these connected complementary clusters. Thus any singular point on $Y$ has a resolution component of weight $<d_1$ above it. By the local isomorphism noted above, the same conclusion applies to a singular $\widetilde x$ of the original model.

A smooth point of an exceptional component of weight $m$, away from its other intersections, has a horizontal point of degree exactly $m$ specializing to it, by Exercise 9.2.11(c); the residue field is infinite and all residue degrees are one. Its image lies in the original formal fibre. Thus singularity in (b) produces a degree $<d_1$ point. At a regular point on the sole component, Exercise 9.2.10 gives degree at least $d_1$, proving the equivalence. In (c), the same inequality at a regular intersection gives degree at least $d_1+d_2$.

\textbf{The minimum bound in (d) is false for unequal multiplicities.} Start with $\mathbf P^1_R$ over $k[[t]]$ in characteristic zero. Blow up a rational fibre point, then the intersection of the resulting two components, then the intersection of the new weight-two component with one weight-one component. The fibre is a chain with weights
\[ 1\;--\;2\;--\;3\;--\;1. \]
The weight-two curve now has square $-2$, and the rightmost weight-one curve has square $-3$. Contract these two disjoint negative rational curves. The resulting normal model has two retained components of weights one and three meeting at an $A_1$ singularity. Its minimal resolution above that intersection has just the weight-two component, so the component demanded in (d) does not exist.

No rational generic point specializes to this intersection. Indeed its lift to the resolution would meet either the weight-two curve alone, giving degree at least two, or one of its intersections with weights one and three, giving degree at least three or five. Properness lifts the horizontal point, so these bounds apply to every point of the original formal fibre. Conversely a smooth point on the weight-two exceptional curve gives degree two. Thus the minimum degree is two, while $\min(1,3)=1$. Locally one can see the same calculation in the normal ring
\[ k[[u,v,w]]/(uv-w^2),\qquad t=uw: \]
the two fibre multiplicities are three and one, and a specialization to the origin requires positive valuations of $u,w$ summing to one, impossible over $k((t))$. This completely disproves both requested conclusions of (d), while (b),(c) retain the stated meanings. \hfill$\square$

\exercise{1.18}{Prove monotonicity and disjoint-subgraph superadditivity of the first Betti number of a graph.}
\pf The cycle space $H_1(G_1,\Q)$ injects into $H_1(G,\Q)$: a graph has no two-dimensional chains, so an already nonzero one-cycle cannot become a boundary after adding edges. This proves (a), or equivalently proves the suggested induction by adding vertices with their incident edges. For pairwise vertex-disjoint subgraphs their cycle spaces have disjoint edge supports, so their direct sum injects into $H_1(G,\Q)$. Taking dimensions proves (b). Disconnected subgraphs are handled componentwise. \hfill$\square$

\exercise{1.19}{Compute the toric rank using the incidence graph of a connected union of connected subcurves.}
\pf Use the reduced curve and geometric branches, as in Exercise 7.5.10. Write $b_x$ for the number of normalization branches at $x$, and $r$ for the number of irreducible components. The normalization-units sequence gives
\[ t(X)=\sum_x(b_x-1)-r+1. \]
For the subcurves use the same formula, replacing $r$ by $r_i$. At a point shared by $m_x$ subcurves, the difference between its contribution for $X$ and the sum of its contributions in the subcurves is $m_x-1$. Since $\sum r_i=r$, subtraction gives
\[ t(X)-\sum_i t(X_i)=\sum_{\text{shared }x}(m_x-1)-n+1. \]
The incidence graph is connected because $X$ and all $X_i$ are connected. It has $n+q$ vertices and $\sum m_x$ edges, where $q$ is the number of shared points. The last displayed expression is exactly $\beta(G)$. \hfill$\square$

\section*{10.2 Reduction of elliptic curves}

\exercise{2.1}{Show that a section of a minimal elliptic model avoids an irreducible component which is a singular conic.}
\pf A section meets the smooth locus of the morphism, as the local retraction makes the base uniformizer a parameter. Over the perfect residue field of Subsection 10.2.1, an integral singular conic consists geometrically of two conjugate lines. Its only rational point, if any, is their singular intersection: a smooth rational point would lie on one line and make it Galois-invariant, contradicting irreducibility. Thus a rational specialization on this component would be a nonsmooth point of the fibre, impossible for a section. \hfill$\square$

\exercise{2.2}{Compute the normal-crossings modifications for types II, III, IV, $\mathrm{IV}_2$, and $\mathrm I_{2n-1,2}$.}
\pf (a) A cusp (type II) starts with multiplicity one. Blow it up: the new curve has multiplicity two and is tangent to the smooth strict transform. Blowing up their tangency creates multiplicity three, leaving a triple point with weights $1,2,3$. Blow up once more to obtain a central curve of multiplicity six, meeting these three disjoint arms transversally. This is Figure 42's type-II picture, with the unmarked strict original curve of multiplicity one. For type III the two multiplicity-one components are tangent. One blow-up creates multiplicity two and a triple point; a second creates multiplicity four. The final star has center four and arms $1,1,2$. Self-intersections at every step follow from Exercise 9.2.9(d), with tangent intersection two decreasing by the product of the local multiplicities.

(b) Type IV has three multiplicity-one components through a triple point. One blow-up inserts a multiplicity-three center and separates the three branches. For $\mathrm{IV}_2$, two geometric branches are conjugate over a quadratic extension. The same blow-up descends: its exceptional curve is $\mathbf P^1_k$ of multiplicity three, one arm is $\mathbf P^1_k$, and the other is $\mathbf P^1_{k'}$ for the quadratic residue extension, meeting the center at a degree-two point. After splitting this is the same three-armed star.

(c) At the nonsplit double point $p$ of the singular conic $\Gamma_n$, its order is two and its fibre multiplicity is one. Blowing up gives $\mathbf P^1_k$ of multiplicity two; the strict transform is its normalization $\mathbf P^1_{k(p)}$ and meets the exceptional curve at one degree-two point. The other components and their multiplicities are unchanged, giving Figure 44.

(d) In the remaining listed types every component is regular and all intersections are transverse with at most two local branches. Thus their divisor already has normal crossings. This assertion uses the perfect-field classification; over imperfect fields a geometrically nonreduced component requires separate care. \hfill$\square$

\exercise{2.3}{Bound by 24 the degree of a separable extension giving good or multiplicative elliptic reduction.}
\pf If the residue characteristic is not two, split the cubic defining the nonzero two-torsion. It is separable and its splitting field has degree at most six. Choose two roots with difference of least valuation; translate and scale so all three roots become $0,1,\lambda$ with $\lambda$ integral. Adjoin the square root needed for the $y$-scale, of degree at most two. Thus a separable extension of degree at most twelve gives a Legendre equation with integral $\lambda$. If $\lambda$ and $1-\lambda$ are units it has good reduction; otherwise its special cubic is nodal and it has multiplicative reduction, possibly nonsplit. No further quadratic splitting extension is required.

In residue characteristic two, $3$ is invertible in $K$. Adjoin one nonzero three-torsion point, at separable degree at most eight. Its rational tangent inflection line gives an equation
\[ y^2+axy+by=x^3. \]
If $\nu(a)\le\nu(b)/3$, scale by $a$ to make $a$ a unit and $b$ integral. Otherwise adjoin a cube root $s$ of $b$ and scale $x=s^2x'$, $y=s^3y'$; now $b'=1$ and $a'$ is integral. This additional extension has separable degree at most three. The normalized equation has $a$ or $b$ a unit. Its discriminant is $b^3(a^3-27b)$. In residue characteristic two the three cases are: a unit discriminant (good); $b$ nonunit and $a$ unit (a node); or $b$ unit and $a^3=27b$ modulo the maximal ideal (again a node after translating its singular point). These are exactly the computations of Proposition 2.33, and give good or multiplicative reduction. The total degree is at most $8\cdot3=24$. \hfill$\square$

\exercise{2.4}{Identify an elliptic curve with its degree-zero Picard functor, compatibly with field extensions.}
\pf If $L$ has degree zero, Riemann--Roch on the genus-one curve gives $h^0(L(o))=1$ and $h^1(L(o))=0$. Its unique effective degree-one divisor is a rational point $x$, so $L=\mathcal O(x-o)$. Uniqueness follows because two effective degree-one divisors in the same class are equal. This proves bijectivity for every extension $K'/K$. The elliptic addition law is defined by the relation $x+y\sim(x+y)_{\mathrm{group}}+o$, so the bijection is a homomorphism. Divisors and their associated invertible sheaves commute with base change, proving functoriality. More generally this is an isomorphism of the represented group functors: the Abel map to $\operatorname{Pic}^0$ is a morphism, and the relative Riemann--Roch construction just used gives its inverse in families. Hence $E\simeq\operatorname{Jac}(E)$ as group schemes, not only as groups of rational points. \hfill$\square$

\exercise{2.5}{Study translations on the smooth locus $N$ and extend the translation family over an imperfect separably closed residue field.}
\pf (a) A section $x$ gives a generic translation. Minimal-model uniqueness extends it to an automorphism $t_x$ of the regular model $\mathcal E$, and it preserves its smooth locus. The rational family $t$ restricts to this translation because the two maps agree generically.

(b) The graphs in $N\times_SN$ are closed Cartier divisors in a regular threefold. Their generic fibres are disjoint when $x\ne x'$. Any nonempty intersection is therefore supported in the special fibre and is of pure codimension two; its first projection identifies each component with a full irreducible component of $N_s$. Such a smooth component has a rational point over the separably closed $k$, and completeness lifts it to a section $z\in N(R)$. Equality at that point would say $t_x(z)_s=t_{x'}(z)_s$. But generic translations commute, so $t_x(z)=t_z(x)$ and $t_{x'}(z)=t_z(x')$ as sections, by separatedness. Since $t_z$ is an automorphism, this contradicts $x_s\ne x'_s$. Thus the graphs are actually disjoint. A possible union of intersection components is treated by choosing any one of them; no irreducibility assertion for that hypothetical union is needed.

(c) We use the established Neron mapping property as an input to obtain the smooth group law on $N$, without any perfectness hypothesis. First prove the asserted covering. For a closed point $y\in N_s$, the parameters $a$ with $y-a\in U_s$ form a nonempty open of the appropriate smooth component after extending to $k(y)$. Since $k$ is separably closed, $k(y)/k$ is purely inseparable; this open descends topologically through that universal homeomorphism. It contains a $k$-point, by smoothness over a separably closed field. Lift it to a section $a\in N(R)$. Then $y\in t_a(U_s)$. Covering all closed points covers $N_s$, since the union is open.

It remains important to extend the family on $\mathcal E\times_SN$, not just on $N\times_SN$. Let $D$ be its maximal domain of definition and $B$ its closed complement. The projection $\mathcal E\times_SN\to N$ is proper. For every codimension-one point $\xi\in N$, Corollary 9.3.30 makes $\mathcal E\times_S\Spec\mathcal O_{N,\xi}$ a minimal regular elliptic model. Its universal translating section therefore extends to an automorphism by minimal-model uniqueness. Consequently $B$ projects to a closed subset missing every generic point of $N_s$, and $D$ contains $\mathcal E\times_SU$ for an open $U$ whose special fibre is dense in $N_s$. The generic fibre is already entirely in $D$.

For any section $a$, second-coordinate translation $\sigma_a(z,x)=(z,x+a)$ is an automorphism of $\mathcal E\times_SN$. The generic identity
\[ t\circ\sigma_a=\alpha_a\circ t,\qquad
\alpha_a(w,x)=(t_a(w),x+a) \]
uses only the individual automorphism $t_a$ from (a), so $\alpha_a$ is also everywhere defined. It follows that $D$ is invariant under all these second-coordinate translations. The covering proved above now gives $D=\mathcal E\times_SN$. Apply the same argument to the rational inverse (translation by the negative parameter). Thus the extended $t$ is an automorphism on the full product, including the nonsmooth points of $\mathcal E$. This proves the desired family extension, with the Neron mapping theorem explicitly declared as the input for the group structure on $N$.

(d) For types $\mathrm I_1$ and II there is a unique singular point on the whole special fibre. Every automorphism, in particular a translation, fixes it. Thus distinct translations can have intersecting graphs in $\mathcal E\times_S\mathcal E$, even though their graphs are disjoint in the smooth locus when the sections have different specializations. \hfill$\square$

\exercise{2.6}{Spread an Abelian variety to an Abelian scheme on a dense open and show that the fibrewise identity components of its Neron model form an open group subscheme.}
\pf Spread a projective embedding of $A$, its smoothness, group law, inverse and identity over an open $V\subset S$. After shrinking, the group identities still hold and all fibres are geometrically connected by openness of geometric connectedness in a proper smooth family. This gives a smooth projective group scheme $A'\to V$.

The Neron mapping property gives $A'\to\mathcal A_V$ extending the generic identity. The inverse generic map spreads to $\mathcal A_U\to A'_U$ after shrinking to an open $U\subset V$: finite presentation and quasi-compactness allow finitely many affine charts, and their maps and equalities spread. Both composites are identities on a further shrinking, proving $\mathcal A_U\simeq A'_U$. Thus the Neron model has geometrically connected fibres over $U$.

Only finitely many fibres remain. In each, the identity component of a smooth algebraic group is geometrically connected and open and closed. The union of the other components is closed in that fibre and hence closed in $\mathcal A$. Removing these finitely many closed subsets gives exactly $\mathcal A^0$. It is open, contains the identity, and is preserved by multiplication and inverse on every fibre, so it is an open group subscheme. \hfill$\square$

\exercise{2.7}{Classify a nonsplit one-dimensional torus and give its norm-one equation.}
\pf The hypothesis that $T$ \emph{is a torus} must be retained: dimension one and $T\not\simeq\mathbf G_m$ alone would allow $\mathbf G_a$, a counterexample to (a). For a torus its character lattice is $\Z$ with a continuous Galois action through $\operatorname{Aut}(\Z)=\{1,-1\}$. A nonsplit action therefore has quadratic separable splitting field $l/k$.

The smooth projective completion is a genus-zero curve with the identity as a rational point, hence $\mathbf P^1_k$. Its two geometric boundary points are interchanged by Galois and descend to one point $q$ with residue field $l$. This proves (a). If $f:T_l\simeq\mathbf G_{m,l}$, the descent discrepancy $\sigma^{-1}f\sigma f^{-1}$ is a group automorphism of $\mathbf G_m$. Such an automorphism is $z\mapsto z^{\pm1}$, since it must act invertibly on the character lattice. It is not the identity, or $f$ would descend. Hence it is inversion. Fixed-point descent now gives
\[ T(k)\simeq\{z\in l^*:z\sigma(z)=1\}. \]
This identifies the group scheme with the norm-one torus as well. If $\theta^2+a\theta+b=0$, write $z=X-\theta Y$. Its norm is $X^2+aXY+bY^2$. Thus the requested affine equation is $X^2+aXY+bY^2=1$, with its norm multiplication law. The argument includes separable quadratic extensions in characteristic two. \hfill$\square$

\exercise{2.8}{Reprove potential good reduction at two for $y^2=x^4-1$.}
\pf The curve has rational points at infinity and hence is elliptic after choosing one as origin. Put $r=1/(x-1)$ and $s=y/(x-1)^2$. Then $s^2=4r^3+6r^2+4r+1$. With $X=4r+2$, $Y=4s$ this is
\[ Y^2=X^3+4X. \]
Its $j$-invariant is $1728$, which belongs to $\Z_2$. Proposition 2.33 therefore gives potential good reduction at two. This proof computes the invariant and uses the potential-good-reduction criterion, rather than constructing the explicit extension in Example 1.28. \hfill$\square$

\exercise{2.9}{Give an elliptic curve with potential good reduction at every finite prime, and rule out potential multiplicative reduction at all primes.}
\pf Take $E:y^2=x^3+1$, for which $j(E)=0$. Its $j$-invariant is integral at every prime, so Proposition 2.33 gives potential good reduction at every prime, including two and three. For any elliptic curve over $\Q$, $j(E)\in\Q$ has negative valuation at only finitely many primes. Potential multiplicative reduction requires a negative valuation of $j$. Thus no elliptic curve over $\Q$ can have it at every prime. \hfill$\square$

\section*{10.3 Stable reduction of algebraic curves}

\exercise{3.1}{Identify arithmetic genus and $h^1$ for a stable curve.}
\pf A stable curve is proper, geometrically reduced and geometrically connected. After algebraic closure its only global functions are constants. Flat base change of $H^0$ gives $H^0(C,\mathcal O_C)=k$. Thus $p_a(C)=1-\chi(\mathcal O_C)=h^1(C,\mathcal O_C)$. \hfill$\square$

\exercise{3.2}{Prove that ample dualizing sheaf characterizes stability among projective geometrically connected semistable curves.}
\pf Extend to an algebraic closure. For every normalized component $\widetilde\Gamma$, let $n_\Gamma$ count the branches above all nodes, including twice each self-node. The residue description gives
\[ \deg_\Gamma\omega_C=2g(\widetilde\Gamma)-2+n_\Gamma. \]
Ampleness makes every such integer positive, so rational normalized components have at least three special branches. Equivalently, a smooth rational component meets the others at least three times; nodal rational components satisfy the same stability count with their self-nodes included. Also $2p_a(C)-2=\deg\omega_C>0$, so $p_a(C)\ge2$. These are the stability conditions. They descend to $k$. \hfill$\square$

\exercise{3.3}{Analyze the natural composite from differentials to the normalization's regular differentials.}
\pf Pullback and the universal derivation give the displayed canonical homomorphisms, and $\Omega^1_{C'/k}=\omega_{C'/k}$ since the normalization of a nodal curve is smooth. \textbf{Their composite is not an isomorphism at a node.} On the split completed node, writing the completed differential module of an algebraic nodal chart,
\[ A=k[[u,v]]/(uv),\qquad
\Omega^1_{A/k}=(A\,du\oplus A\,dv)/(v\,du+u\,dv), \]
the nonzero element $u\,dv=-v\,du$ is killed by both $u$ and $v$ and maps to zero on the normalization. Conversely $du$ and $dv$ generate $k[[u]]du\oplus k[[v]]dv$, so the composite is surjective, with a one-dimensional torsion kernel at this node. Etale descent gives the same assertion at nonsplit nodes, with residue degree included. The corrected isomorphism is
\[ \Omega^1_{C/k}/\text{torsion}\ \simeq\ \pi_*\omega_{C'/k}. \]
The dualizing sheaf $\omega_C$ itself is larger, allowing simple poles of opposite residues on the two branches. \hfill$\square$

\exercise{3.4}{Bound the normalization fibre over a point of a semistable curve.}
\pf At a smooth point the normalization is an isomorphism and the dimension is one. At a split node its completed normalization is $k(x)[[u]]\oplus k(x)[[v]]$ and reduction modulo the maximal ideal is $k(x)\oplus k(x)$, of dimension two. At a nonsplit node this algebra is a separable quadratic field instead, still of dimension two. Completion is faithfully flat and preserves this finite fibre, so the asserted bound holds before completion. In particular there are at most two points over $x$. \hfill$\square$

\exercise{3.5}{Characterize semistability using the normalization quotient, residue separability and the normalization fibre.}
\pf At a node, $k(x)/k$ is separable, the normalization fibre is a finite \emph{etale} algebra of dimension two, and the normalization quotient has length one over $\mathcal O_{C,x}$. These conditions also imply a node: length one makes the maximal ideal annihilate the quotient; the conductor description then identifies the completed local ring with the two-branch fibre product, or its nonsplit separable quadratic form. This is the usual normalization criterion for an ordinary double point.

\textbf{Reduced must be strengthened to geometrically reduced in the printed fibre condition.} In characteristic two choose $a\in k\setminus k^2$ and
\[ C=\Spec k[x,y]/(y^2-x^3-ax^2). \]
Its sole nonsmooth point is $(0,0)$ with residue field $k$. Its normalization is $k[T]$, with $x=T^2-a$, $y=T(T^2-a)$, and its ring is $k+(T^2-a)k[T]$. Thus the normalization quotient has dimension one. The normalization fibre is $k[T]/(T^2-a)$, a reduced purely inseparable quadratic field. All the printed conditions hold, but after adjoining $\sqrt a$ the equation becomes $(y-\sqrt a\,x)^2=x^3$, a cusp, not a node. Requiring the finite fibre geometrically reduced, hence etale, excludes exactly this failure. The normalization quotient is interpreted as length one, or as a one-dimensional $k(x)$-space once its maximal ideal annihilation is known. \hfill$\square$

\exercise{3.6}{Describe automorphisms via the normalization and test finiteness as a criterion for stability.}
\pf (a) Every automorphism lifts uniquely to the normalization and permutes the unordered pairs of points lying above nodes. Conversely an automorphism preserving this \emph{collection} of pairs descends through the nodal gluing, by Proposition 3.7(c). It need not fix each pair individually, nor each smooth point. Thus the printed condition should say that the normalization equivalence relation is preserved, not $\sigma(\pi^{-1}(x))=\pi^{-1}(x)$ for every individual $x$. For example an automorphism may interchange two nodes.

(b) Over an algebraically closed field the intended criterion is correct. On each normalized component, stability gives either genus at least two, genus one with a marked point, or genus zero with at least three marked points. Their automorphism groups preserving these finite sets are finite, and there are only finitely many component and node permutations. Conversely an unstable rational component has at most two marked branches and an infinite subgroup of $\operatorname{PGL}_2$ fixing them; extend it by the identity on the other components. A smooth genus-one curve has its infinite translation group. These give the converse.

\textbf{Over merely infinite $k$, genus one is a counterexample.} Take $k=\mathbf F_5(t)$ and the constant elliptic curve $E:y^2=x^3-x$ over $k$. Every $k$-point defines a rational map $\mathbf P^1_{\mathbf F_5}\dashrightarrow E$, which extends by properness. It is constant: a nonconstant map from a genus-zero curve to a genus-one curve is impossible by Hurwitz after separating off Frobenius. Thus $E(k)=E(\mathbf F_5)$ is finite. Its pointed automorphism group is also finite, so its full group $E(k)\rtimes\operatorname{Aut}_k(E,o)$ is finite, although this smooth genus-one curve is not stable. With the additional assumption $p_a(C)\ge2$ and split nodes, the intended proof works over any infinite $k$: an unstable rational component has rational attachment points and is $\mathbf P^1_k$, whose subgroup fixing at most two points is infinite. The exceptional unmarked genus-one case is then absent. \hfill$\square$

\exercise{3.7}{Test the claim that the canonical sheaf of a stable curve is generated exactly when the curve is geometrically integral.}
\pf \textbf{The converse direction is false.} Glue two smooth elliptic curves along two distinct pairs of rational points. This is a stable reducible curve of genus three. On each component the dualizing sheaf is $\omega_E(p+q)=\mathcal O_E(p+q)$. A nonzero holomorphic differential, extended by zero on the other component, generates at every smooth point of that component. At the nodes choose differentials of the third kind on both components with residues $1,-1$, and opposite residues on the paired branches. They glue to a global dualizing section nonzero at both nodes. Thus $\omega_C$ is globally generated despite reducibility.

The correct geometric criterion is the absence of \emph{separating nodes}. A global dualizing section is a differential on the normalization with simple poles whose residues cancel across each node and sum to zero on each normalized component. Residues therefore form a flow on the dual graph. A bridge has zero residue for every flow, forcing every section to vanish at its node. A nonbridge lies on a cycle, giving a flow nonzero there and hence a section generating at that node. At a smooth point of a positive-genus component use holomorphic differentials. On a rational component choose a cycle through its vertex; the corresponding third-kind differential has exactly two poles there and no smooth zeros. This generates at its smooth points as well. These arguments prove the criterion over an algebraic closure, and global generation descends by base change. An integral stable curve has no separating node, proving the valid forward assertion. \hfill$\square$

\exercise{3.8}{Characterize semistable conic reduction by etale-local splitting and examine the proposed finite Galois refinement.}
\pf (a) A geometric semistable genus-zero fibre is a tree of smooth rational curves. It has a smooth separable point. After an etale residue extension that point is rational, and the smooth local-section lemma lifts it to a section over an etale DVR extension. Its generic point splits the conic.

(b) The finite Galois refinement in (1) is valid over a Henselian $R$, but false for a general DVR when $R_L/R$ is only etale, not finite. For an explicit obstruction let $K=\Q$, $R=\Z_{(3)}$, and let $L=\Q(\alpha)$ for $P=X^3-X^2+3$. This polynomial is irreducible by the rational-root test, has reduction $X^2(X-1)$, and has discriminant $-231$ of three-adic valuation one. Its simple branch at one is unramified; dividing by its Hensel root leaves an Eisenstein quadratic factor, so the branch at zero is ramified. Localize its integral closure at the simple branch: this is etale over $R$, but the field has another ramified place above three. It cannot be contained in the fraction field of a finite etale Galois $R$-algebra, since every place of every subfield of such an algebra is unramified. Taking $C=\mathbf P^1_K$ satisfies the conic-splitting hypothesis as well, so this disproves (1) in its stated context without the Henselian hypothesis.

Here is a proof of the desired semistability that avoids that erroneous global refinement. Let $X/R$ be relatively minimal. Its fibre is integral by Exercise 9.3.1. Pass to a strict Henselization where the conic splits by the hypothesis. The regular base change is then a regular model of $\mathbf P^1$, so its geometric fibre graph is a tree. Galois acts transitively on its components because the original fibre is irreducible. A finite tree with a vertex-transitive automorphism group has one vertex or two adjacent vertices: its center is a fixed vertex or fixed edge. Moreover all multiplicities are equal by conjugacy and one is one because the split generic curve has a section. All are therefore one. The components are smooth projective lines, and the only possible intersection is an ordinary node. Thus the base-changed fibre, and by etale descent the original fibre, is semistable. Over a Henselian base the finite Galois cover in (1) exists and gives exactly the action and transitivity requested in (2). Together with (a) this proves the correct conic criterion, without confusing it with good reduction. \hfill$\square$

\exercise{3.9}{Show that semistability is invariant under surjective etale morphisms on the source.}
\pf Flatness of $X\to S$ is equivalent to flatness of $Y\to S$ by faithfully flat descent along $X\to Y$; the converse uses composition with an etale, hence flat, map. On every geometric fibre, the map is again etale and surjective. Smooth points and ordinary nodes are exactly the etale-local models $\mathbf A^1$ and $uv=0$, and this property ascends and descends etale locally. Thus all geometric fibres are reduced nodal curves on one side exactly when they are on the other. Finite type was assumed, so this proves both implications. Properness is not part of this semistability assertion. \hfill$\square$

\exercise{3.10}{Analyze the stable quartic $((x^2-2y^2+z^2)(x^2-z^2)+p^ry^3z)=0$ and compute its node thicknesses.}
\pf Put $c=p^r$, with $p>2$. The model is flat and projective. Its special fibre is the smooth conic $x^2-2y^2+z^2=0$ and the two lines $x=z$, $x=-z$. The lines meet at $q=[0:1:0]$, away from the conic. Each line meets the conic at two distinct rational points, giving four more transverse nodes. Each rational line has three marked points and the conic has four, so the fibre is stable. The generic quartic is smooth: in the chart $y=1$, the $x$ derivative forces a possible singular point to have $x=0$ or $x^2=1$. Combining the equation and the $z$ derivative would then force respectively $c^2=32/27$ or $c^2=-256/27$, impossible for the rational number $c=p^r$. Points with $y=0$ are smooth by the derivatives of the two quadratic factors. Thus this is a stable family.

At any of the four conic-line nodes $y,z$ are units, so the first-order smoothing term is $c$ times a unit; its thickness is $r$. At $q$, set $y=1$. The equation is
\[ x^4-z^4-2x^2+2z^2+cz=0. \]
Its nondegenerate quadratic part permits the formal Morse lemma, since two is invertible. Its critical point has $x=0$ and $z=-c/4+O(c^3)$; substitution gives critical value $-c^2/8+O(c^4)$. Hence its completed local equation is $UV=\text{unit}\cdot p^{2r}$, proving thickness $2r$. The other points have completed equation $UV=\text{unit}\cdot p^r$. \hfill$\square$

\exercise{3.11}{Find a quadratic Dedekind base change making the minimal desingularization of a semistable family a normal-crossings regular model.}
\pf For the term arithmetic surface assume $X$ integral; a disconnected union cannot become integral after base change, so without this assumption the conclusion is componentwise. Only finitely many fibres have nodes. Include one auxiliary closed base point among the prescribed ramification points if necessary. Choose a separable quadratic extension of $K(S)$ ramified at each of their base points. If $\operatorname{char}K\ne2$, weak approximation supplies $a$ of valuation one there, and take $K(\sqrt a)$. In characteristic two choose $a$ of valuation $-1$ there and take $w^2-w=a$. The odd pole forces ramification index two. The normalization $S'$ is finite flat of degree two and Dedekind. The generic constant field of a proper semistable integral family extends to a finite etale Stein base. Ramification at the auxiliary point makes the quadratic extension disjoint from that constant field, so the base-changed generic curve stays integral. Finiteness for this separable extension also follows directly by bounding integral elements by the trace-dual of a finite integral basis.

After base change a node of thickness $n$ has thickness $2n$. Its minimal resolution inserts a chain of $2n-1$ rational curves of multiplicity one. In particular at least one curve is inserted, separating the two branches even if the original component had a self-node. All strict components and inserted curves are regular and meet transversally, with no triple points. Smooth parts stay smooth. Hence the minimal desingularization is a regular model with normal crossings. The construction descends from splitting etale neighborhoods of nonsplit nodes, so a split-node assumption is unnecessary. \hfill$\square$

\exercise{3.12}{Show that every proper semistable model of a smooth genus-at-least-two curve dominates its stable model.}
\pf Resolve the surface nodes; the exceptional fibres are rational chains. The resulting regular semistable model dominates the minimal regular model. Contract its rational tails and then its rational bridges to obtain the stable model, as in Theorem 3.34. Every exceptional chain introduced by the initial resolution is contracted in this stabilization: it has genus zero and at most two attachments. Therefore the map to the stable model is constant on every fibre of the resolution. The original semistable model is normal, and its resolution pushes its structure sheaf down to that of the model, so the contraction descent criterion gives a morphism from the original model. It is unique since it is the identity on the generic fibre. If a ``model'' is not required proper, the assertion must include that hypothesis; an arbitrary open model need not dominate a proper one by a surjective birational morphism. \hfill$\square$

\exercise{3.13}{Descend semistable reduction from a smooth surjective base change.}
\pf \textbf{Properness of the family $X\to T$ is missing.} Let $S=T=\Spec k[[t]]$ in characteristic zero, and take the smooth locus of the regular cubic model $y^2z=x^3+tz^3$. Its generic fibre is the complete elliptic curve, and it is smooth, hence semistable, over $S$. But its generic elliptic curve has additive type-II reduction and no proper semistable model over $R$. Thus a nonproper $X$ as permitted by Definition 3.14 does not imply the claim.

For the corrected assertion assume $X\to T$ proper. At each closed $s\in S$, the smooth surjection $T\to S$ has a section after passage to an etale neighborhood, by Exercise 6.2.7(b). Pulling back $X$ gives a proper semistable model of the corresponding generic curve over an etale DVR extension of $\mathcal O_{S,s}$. Semistable reduction descends along these etale extensions, by Corollary 3.36 (for conics use 3.8). Therefore $C$ has semistable reduction at every $s$. The local proper models glue through the minimal regular or stable model, giving reduction over $S$. The conclusion concerns $C$, rather than the printed ``$X$ over $S$''. \hfill$\square$

\exercise{3.14}{Construct a common etale local neighborhood of a split node and $k[u,v]/(uv)$.}
\pf (a) If the local ring $A_1$ is not integral, the two branches are smooth and their ideals are generated by parameters $u,v$ with $uv=0$. The map from $A_0$ defined by these parameters induces an isomorphism on completions, and hence is etale by the local flatness and unramified criteria. Thus $R=A_1$ works.

(b) For integral $A_1$, its semilocal normalization $B$ has two rational closed points. The Chinese remainder theorem gives $f\in B$ with values one and zero. The conductor criterion for a node says an element of $B$ lies in $A_1$ exactly when these values agree. Hence $a=f^2-f\in A_1$, and $a$ is in its maximal ideal.

(c) The algebra $D=A_1[w]/(w^2-w-a)$ is finite etale: its discriminant $1+4a$ is a unit, also in characteristic two. If $f=a_1/a_2$, direct multiplication gives
\[ (a_2w-a_1)(a_2w-a_2+a_1)=0. \]
Generically the two factors correspond to $f$ and $1-f$. At either closed point of $D$, both generic branches meet: their residues on the two normalization branches are interchanged, by the chosen values of $f$. Thus the local ring $D_\fm$ is nonintegral, and is a split node because it is etale over $A_1$ with rational residue field. Apply (a) to it. The composite maps give the common etale local ring over $A_0$ and $A_1$. \hfill$\square$

\exercise{3.15}{Construct a common etale local neighborhood of a split thickness-$n$ node and $uv=t^n$.}
\pf (a) Suppose first the special local ring has two distinct branches. The total local ring $B_1$ is normal, with completed ring $\widehat R[[U,V]]/(UV-t^n)$. Let $\Gamma_1,\Gamma_2$ be its two vertical prime divisors. In the completion, $n\Gamma_1$ and $n\Gamma_2$ are principal. Their divisorial ideals are finite reflexive modules; duals commute with flat completion, so these ideals become free after completion and were already invertible over the local ring by faithful flatness. Choose generators $u,v$ with divisors $n\Gamma_1,n\Gamma_2$. Their product has divisor $nX_s$, hence is a unit times $t^n$. Absorb that unit in $v$. Their reductions generate the two branch parameters, and the resulting map from $B_0$ induces an isomorphism of completions, hence is etale. Thus $B_1$ itself is a common neighborhood.

The special generators must be allowed this unit normalization. Arbitrarily prescribed generators need not lift with product exactly $t^n$. Already in $uv=t$, prescribing $2\overline u,\overline v$ over a characteristic-zero residue field makes the product's coefficient of $t$ equal to two modulo the maximal ideal, which cannot be changed to one by adding multiples of $t$. This is a harmless normalization correction to the printed wording.

(b) In the integral special-node case choose $a$ from 3.14 and lift it to $b\in B_1$. Since $b$ belongs to the maximal ideal, $B_1[w]/(w^2-w-b)$ is finite etale. At a closed point its special fibre is the nonintegral split node constructed in 3.14. Apply (a) to that localization. Composing etale maps gives the requested common neighborhood. No general approximation theorem is needed for this node calculation. \hfill$\square$

\exercise{3.16}{Replace the reduced semistable special pieces in Proposition 3.44 by smooth reductions, allowing multiplicities, and show that the stable graph is a tree.}
\pf We give a proof using the Jacobian specialization and semistable-reduction criteria, explicitly stronger inputs than the elementary reduced-piece proof of Proposition 3.44. These say: the Abelian rank of the Neron special fibre of $J=\operatorname{Jac}(C)$ equals the sum of genera of the normalized components of any regular model; and, over a perfect residue field, a smooth curve of genus at least two has stable reduction if and only if its Jacobian has semi-Abelian reduction (Deligne--Mumford, Theorem 2.4). Write $\mathcal J$ for its Neron model. In particular an Abelian rank equal to $\dim J$ means good reduction of $J$.

Pass to a complete strictly Henselian base, which does not change the question, and take a common regular model of the finitely many pieces. Their distinct component valuations give distinct strict components. Because their reduced special curves $U_i$ are smooth, their normalized projective completions contribute the sum of their genera, which is at least $p_a(\widehat U_i)$ (with equality for connected $U_i$), to the Abelian rank. Consequently
\[ a(\mathcal J_s)\ge\sum_i p_a(\widehat U_i)\ge g(C)=\dim J. \]
The opposite bound is automatic, so equality holds, and $J$ has good reduction. The cited criterion supplies the stable model of $C$ over the original DVR (descent gives the same conclusion before strict completion). Its special Jacobian has Abelian rank $g$ and toric rank zero. The nodal genus formula then gives $\beta(G)=g-\sum g(\widetilde\Gamma)=0$. Thus the connected dual graph is a tree. This proof allows the original pieces to have arbitrary multiplicities, and does not silently apply the reduced-fibre cohomology calculation to a multiple fibre. \hfill$\square$

\exercise{3.17}{Compare Abelian and toric ranks under contraction, and resolve normal models with reduced ordinary-multiple-point special fibre.}
\pf (a) Abelian rank is the sum of the genera of normalized irreducible components. A birational contraction removes exactly those in $E$, so $a(X_s)=a(Y_s)+a(E)$. For toric rank use $t=\sum_x(b_x-1)-r+1$. Collapsing the connected incidence subgraph of $E$ removes its internal cycle space and changes no cycles outside it. Thus $t(X_s)=t(Y_s)+t(E)$. This is also the direct branch-count computation in Lemma 7.5.18(b), and does not require the special fibres reduced as schemes, since these ranks depend on their reductions.

(b) An ordinary-multiple-point reduced fibre has no unipotent Picard part, so $a(Y_s)+t(Y_s)=h^1(Y_s,\mathcal O)=g$. Here $H^0=k$ and flatness gives the Euler characteristic. Cohomological flatness of $X$ gives $h^1(X_s,\mathcal O)=g$. By (a), $a(X_s)+t(X_s)\ge g$, and the identity $h^1=a+t+u$ forces $u(X_s)=0$ and $a(E)=t(E)=0$.

(c) Take the minimal desingularization $X\to Y$. Every strict component has multiplicity one because $Y_s$ is reduced, so the gcd is one and $X$ is cohomologically flat. Part (b) applies. Proposition 3.42 makes its reduced fibre nodal. If the full fibre were nonreduced, the last-step argument in that proposition would produce a multiplicity-at-least-two exceptional $(-1)$-curve. It is exceptional over $Y$, since all retained components have multiplicity one, contradicting minimality of the desingularization. Hence $X_s$ is reduced nodal. Semistability is open in this finite-type flat family; over the local base its open neighborhood of the whole closed fibre also contains the generic fibre. Since that generic fibre is normal, a nodal generic fibre has no singularities and is smooth. Thus $Y_K=X_K$ is smooth and the minimal desingularization is semistable, as claimed. \hfill$\square$

\exercise{3.18}{Analyze decomposition and inertia groups on components of the quotient of an integral semistable family.}
\pf Let $A=\mathcal O_{Y,\xi}$ and let $B$ be its integral closure in $K(X)$. The fibre is reduced on both sides: if an invariant has a power divisible by $t$, reducedness upstairs makes it divisible by $t$ upstairs, and its quotient by $t$ is still invariant. Thus $A$ and the component local rings of $B$ are DVRs with $t$ of valuation one. All ramification indices are one. The fundamental degree formula gives
\[ |G|=\sum_i[k(\xi_i):k(\xi)]=r[k(\xi_1):k(\xi)]. \]
Transitivity on the primes gives $|G|=r|D|$. The standard residue-field theorem for a finite Galois extension of valuation rings identifies the separable residue extension's group with $D/I$; its remaining degree is purely inseparable.

The finite map $X_s/G\to Y_s$ is surjective and universally injective on points. On residue function fields the discrepancy is purely inseparable by the preceding calculation, so the map is radicial. If $|G|$ is invertible in $k$, averaging makes taking invariants exact and commute with reduction modulo $t$. Hence this finite map is actually an isomorphism, not just birational. Also $I$ is then trivial, since its order is the purely inseparable residue degree and must be a power of the residue characteristic. Thus $D$ acts faithfully on the function field, and on the component $\overline{\{\xi_1\}}$. This last action is on the component itself; an action on the finite set of component indices alone would not in general be faithful. \hfill$\square$

\exercise{3.19}{Show that relative Frobenius preserves semistable models and extends a finite purely inseparable map of smooth curves.}
\pf (a) The twist $\mathcal C^{(p)}$ is the base change by absolute Frobenius of $\Spec R$. Semistability is preserved by arbitrary base change, so it is semistable. Relative Frobenius is finite for these finite-type relative curves and purely inseparable.

(b) A separably generated one-variable function field $F/K$ satisfies $[F:KF^p]=p$. If $F/E$ is finite purely inseparable and $E/K$ is also the function field of a smooth curve, successive applications give $E=KF^{p^r}$ for some $r$. Thus the morphism is a sequence of relative Frobenius maps, up to the resulting isomorphism of the target with $C^{(p^r)}$. Indeed the differential of an inseparable map of smooth curves is zero, so $E$ lies in the kernel $KF^p$ of $d:F\to\Omega^1_{F/K}$. Remove the first relative Frobenius step, which has degree $p$, and induct. Define $\mathcal D=\mathcal C^{(p^r)}$. This is a semistable model and the composite relative Frobenius is the required finite purely inseparable extension of $f$. Properness, smoothness and stability, when present, persist under these base changes. \hfill$\square$

\exercise{3.20}{Descend stable reduction against a purely inseparable map of smooth genus-at-least-two curves in equal characteristic, assuming a rational point on the source.}
\pf Pass to the complete base with algebraically closed residue field as allowed in the question. Normalize the stable model $\mathcal D$ in $K(C)$. Completeness gives finiteness of normalization; the extension is purely inseparable, so $\mathcal C\to\mathcal D$ is finite radicial. On reductions it preserves the branch incidence graph and induces purely inseparable maps between normalized component curves. Over the perfect algebraically closed $k$, such maps are iterated Frobenius twists, so their genera are equal. Hence
\[ a(\mathcal C_s)=a(\mathcal D_s),\qquad
 t(\mathcal C_s)=t(\mathcal D_s). \]
Take a regular desingularization $X\to\mathcal C$. By 3.17(a), applied separately to each exceptional fibre, $a(X_s)+t(X_s)\ge a(\mathcal D_s)+t(\mathcal D_s)=g$. The genera of $C,D$ agree by 3.19. A $K$-point of $C$ extends to $X$, forcing the multiplicity gcd to be one. Therefore $X$ is cohomologically flat and $h^1(X_s,\mathcal O)=g$. It follows that its unipotent rank is zero. Proposition 3.42 and stabilization now give a stable model of $C$ over $R$. This proof explains the role of the rational point: it supplies the fibre cohomology equality, rather than merely a point on the normalization. \hfill$\square$

\section*{10.4 Deligne--Mumford theorem}

\exercise{4.1}{Compute the normalized base change of Proposition 4.6 for elliptic Kodaira fibres when the residue characteristic is different from two, three and five.}
\pf We retain the proposition's algebraically closed residue field and choose the minimal regular normal-crossings model. The following recipe specifies every component, attachment and thickness of $\mathcal C'$, including components that are later contracted. Let a component of multiplicity $d$ meet neighbors of multiplicities $a_1,\ldots,a_r$. Its inverse image is the Kummer covering determined by the local unit in $t=u^d b$. On a rational component the number of connected covers is
\[ \ell=\gcd(d,a_1,\ldots,a_r), \]
each of degree $d/\ell$, ramified at the marked points with indices $d/\gcd(d,a_j)$. Indeed the Kummer divisor has these orders modulo $d$, and $\mathbf P^1$ has no nontrivial unramified connected covers. Above an intersection of weights $a,b$ there are $\gcd(a,b)$ nodes, each of thickness $e/\operatorname{lcm}(a,b)$. These assertions are the explicit local integral-closure formulas (4.32),(4.34) and Remark 4.8. They also identify which lifted vertices attach, by their Kummer root labels.

For potential-good types the input graph is a star with one marked central weight $d$ and the following arms; an arm lists weights from center outward:
\[
\begin{array}{c|c|l|c}
\text{type}&d&\text{arms}&\operatorname{lcm}\text{ of all weights}\\ \hline
\mathrm{II}&6&(1),(2),(3)&6\\
\mathrm{III}&4&(1),(1),(2)&4\\
\mathrm{IV}&3&(1),(1),(1)&3\\
\mathrm I_0^*&2&(1),(1),(1),(1)&2\\
\mathrm{IV}^*&3&(2,1),(2,1),(2,1)&6\\
\mathrm{III}^*&4&(3,2,1),(3,2,1),(2)&12\\
\mathrm{II}^*&6&(5,4,3,2,1),(4,2),(3)&60
\end{array}
\]
For explicit equations of those central covers, after placing three marked points at $0,1,\infty$, take the smooth projective normalizations of
\[
\begin{array}{c|l}
\mathrm{II},\mathrm{II}^*&Y^6=X(X-1)^2\\
\mathrm{III},\mathrm{III}^*&Y^4=X(X-1)\\
\mathrm{IV},\mathrm{IV}^*&Y^3=X(X-1)\\
\mathrm I_0^*&Y^2=X(X-1)(X-\lambda),\quad\lambda\ne0,1.
\end{array}
\]
For the starred types invert the Kummer generator and multiply by $X(X-1)$ to replace exponents by their negatives modulo $d$. In the last row $\lambda$ is the cross-ratio of the four attachment points on the original central component. Thus its value is part of the initial elliptic model, not an extra arbitrary choice.

The centers lift to one smooth elliptic component, of cover degree $d$. Hurwitz gives genus one from their respective ramification signatures $(6,3,2)$, $(4,4,2)$, $(3,3,3)$ or $(2,2,2,2)$. Every other lifted component is rational. Explicitly the rational trees attached to the elliptic component are: six single tails for II; four for III; three for IV; four for $\mathrm I_0^*$; three chains of length two for $\mathrm{IV}^*$; two chains of length three and two single tails for $\mathrm{III}^*$; one chain of length five, two chains of length two and three single tails for $\mathrm{II}^*$. All have multiplicity one. Each edge's thickness is given by its original pair of weights, as above. This completely describes $\mathcal C'$ for every admissible $e$. In particular the proposition's divisibility requirement is stronger than the least extension giving a smooth minimal elliptic model: for II* it requires $60\mid e$, not just $6\mid e$.

For $\mathrm I_n$, $n\ge2$, the graph is the original $n$-cycle, all weights one, so $\mathcal C'$ has the same rational components and node thickness $e$. For $\mathrm I_1$ first blow up its self-node: the normal-crossings graph has two vertices of weights one and two joined twice. Its normalization after an admissible even $e$ gives a two-cycle, with thickness $e/2$ at each node. For $\mathrm I_n^*$ with $n\ge1$, the weight-two chain has $n+1$ vertices and four weight-one terminal arms. Its two end vertices lift once, the $n-1$ interior vertices twice. The result is a $2n$-cycle of rational curves with four single rational tails; every node has thickness $e/2$. For good type $\mathrm I_0$, base change is already smooth. Resolving these node thicknesses and contracting rational tails gives respectively good elliptic reduction or multiplicative type $\mathrm I_{en}$. These last contractions describe the minimal model; the preceding lists describe the requested uncontracted normalized model. \hfill$\square$

\exercise{4.2}{Descend stable reduction from a purely inseparable constant extension in equal characteristic.}
\pf Reduce as in the question to a complete $K$ with perfect algebraically closed $k$ and $K'=K^{1/p}$. Frobenius on $R'$ factors through $R$, because every $p$th power in $R'$ lies in $R$ (valuation integrality and $K'^p=K$). Let $\mathcal C'/R'$ be the stable model. Its relative Frobenius twist is the base change of an $R$-scheme $\mathcal D$: explicitly base change the same equations of $\mathcal C'$ along $R'\to R$, $a\mapsto a^p$, and then along $R\to R'$. This yields $\mathcal C'^{(p)}=\mathcal D_{R'}$.

Stable curves descend under the faithfully flat extension $R'/R$: properness, flatness, nodal fibres and ample relative dualizing sheaf all descend. Its generic fibre descends with the displayed Frobenius-twisted equations and is $C^{(p)}$. The purely inseparable map $C\to C^{(p)}$ and the rational-point hypothesis now satisfy 3.20, which gives stable reduction of $C$. Iterate for higher purely inseparable degree. This argument does not try to descend an arbitrary model directly along a field extension lacking Galois automorphisms. \hfill$\square$

\exercise{4.3}{Algebraize a finite separable extension of the completed fraction field by sufficiently close polynomial approximation.}
\pf Choose an integral primitive element $\pi$ with monic minimal polynomial $P$ of degree $n$, and write $c=P'(\pi)\ne0$. Interpret coefficient divisibility by $tc^2$ inside $R_L$; it suffices to approximate all coefficients to an arbitrarily large $t$-adic order. Choose a monic separable $Q\in R[X]$ of degree $n$ with
\[ \nu_L(Q_i-P_i)\ge\nu_L(t)+2\nu_L(c). \]
This is possible by density of $K$ in $\widehat K$; separability follows as well from the estimates below.

(a) If $\varepsilon\in tcR_L$, Taylor expansion gives $P'(\pi+\varepsilon)-c\in tcR_L$, since $P$ has integral coefficients. The derivative approximation has still larger valuation. Thus $\nu_L(Q'(\pi+\varepsilon))=\nu_L(c)$.

(b) Newton iteration starting at $\pi$ has first correction of valuation at least $\nu_L(t)+\nu_L(c)$, and subsequent errors improve quadratically. Completeness gives a root $y\in R_L$ with
\[ \nu_L(y-\pi)\ge\nu_L(t)+\nu_L(c)>\nu_L(c). \]
In particular $y-\pi\in\fm_L$.

(c) All conjugates of the integral $\pi$ are integral and $c=\prod_{\pi'\ne\pi}(\pi-\pi')$. Therefore each conjugate distance has valuation at most $\nu_L(c)$. The stronger bound in (b) and Krasner's lemma give $\widehat K(\pi)\subseteq\widehat K(y)\subseteq L$. Hence equality holds. Since $Q$ has degree $n=[L:\widehat K]$, it is irreducible over $\widehat K$, hence over $K$. Thus $K'=K[X]/(Q)$ is a finite separable field and $K'\otimes_K\widehat K=L$. Mere congruence $y\equiv\pi$ would not alone prove this generation claim; the conjugate-distance estimate is essential.

(d) The integral closure $R'$ of $R$ in this separable field is finite, by the trace-dual bound. Its completed base change is the product of the completed DVRs at its maximal ideals. Their fraction algebras form $K'\otimes_K\widehat K=L$, a single field, so there is only one such factor. It is the complete valuation ring of $L$, namely $R_L$. Therefore $\widehat R\otimes_RR'\simeq R_L$. \hfill$\square$

\exercise{4.4}{Algebraize a compatible inverse system of coherent sheaves on the thickenings of a projective scheme over a complete DVR, and apply this to the Picard group.}
\pf (a) Compatibility gives $F_N=F_{N+n}/t^{N+1}F_{N+n}$ and the first asserted exact sequence. Multiplication by $t^N$ gives a well-defined surjection
\[ \rho_{N,n}:t^nF_n\longrightarrow t^{N+n}F_{N+n}. \]
In particular the kernels of $F_0\twoheadrightarrow t^NF_N$ increase. Noetherianity of the coherent $F_0$ makes them stationary. If $N_0$ is a stationary index and $G$ the stabilized quotient, $t^NF_N\simeq G$ for all $N\ge N_0$.

(b) Choose $d_1$ so that $F_{N_0}(d_1)$ is generated and $H^1(G(d_1))=0$. The exact sequences
\[ 0\to G(d_1)\to F_{N+1}(d_1)\to F_N(d_1)\to0 \]
for $N\ge N_0$ lift all global sections. Lift a fixed finite set of generators compatibly through the system. Nakayama on each nilpotent thickening shows the lifted sections generate every $F_N(d_1)$. They define compatible surjections $L_{1,N}\twoheadrightarrow F_N$ with $L_1=\mathcal O_X(-d_1)^r$.

(c) On each completed affine chart the inverse limit is a finite module over the complete Noetherian coordinate ring, since the lifted generators just constructed generate modulo every $t^{N+1}$. Its kernel in $\widehat L_1$ is finite. These kernels glue to a coherent formal sheaf $\mathcal K$; their reductions modulo $t^{N+1}$ are compatible coherent sheaves on $X_N$. This kernel construction is also obtained by Artin--Rees from the stabilized images of the ordinary kernels. One must not assume the raw kernels $\ker(L_{1,N}\to F_N)$ commute with reduction when there is torsion.

Apply (b) to the reductions of $\mathcal K$. It supplies a sum of negative ample twists $L_2$ and a compatible formal morphism $\widehat L_2\to\widehat L_1$ with cokernel the original system. Formal functions, using that $R$ is complete and $L_2$ locally free, gives
\[ H^0(X,\mathcal Hom(L_2,L_1))
\simeq\varprojlim_N H^0(X_N,\mathcal Hom(L_{2,N},L_{1,N})). \]
Thus this morphism algebraizes to $f:L_2\to L_1$. With $F=\operatorname{coker}f$, right exactness gives $F|X_N=F_N$. The printed direction $L_1\to L_2$ and $\mathcal Hom(L_1,L_2)$ must both be reversed.

(d) If all $F_N$ are invertible, at a closed-fibre point choose a lift of a basis modulo $t$. It generates $F$ by Nakayama and is a basis modulo every $t^{N+1}$. Its kernel lies in every power of the maximal ideal and is zero by Krull intersection. Thus $F$ is locally free of rank one near the closed fibre. The non-locally-free locus is closed and proper over the local base; if nonempty its image would contain the closed point. It is therefore empty. Finally choose compatible representatives and restriction isomorphisms for any element of $\varprojlim\operatorname{Pic}(X_N)$ and apply this algebraization. It gives the asserted surjectivity of $\operatorname{Pic}(X)$. \hfill$\square$

\exercise{4.5}{Relate the fibre multiplicity gcd to the residue characteristic when the degree-zero Picard group is unipotent.}
\pf First assume $H^0(X_\eta,\mathcal O)=K$, as holds for a geometrically connected generic curve. If a prime $\ell\ne\operatorname{char}k$ divides $d$, the divisor $X_s/\ell$ is integral and its line bundle $L$ has exact order $\ell$: it has $\ell$th power trivial, but if it were principal, its rational function would be a generic constant in $K^*$, whose vertical divisor is an integral multiple of the full fibre. Its restriction to $X_s$ has multidegree zero, so lies in the unipotent group $\operatorname{Pic}^0(X_s)$. Such a group over algebraically closed $k$ has no prime-to-characteristic torsion. Thus $L|X_s$ is trivial. Lemma 4.14(a) forces its torsion restriction to every thickening to be trivial, and part (b) then forces $L$ trivial on $X$, a contradiction. Hence $d$ is a power of the residue characteristic; in characteristic zero this means $d=1$.

\textbf{Without the generic-constant-field hypothesis the printed assertion is false.} Take $R=k[[t]]$, $R'=k[[u]]$, $t=u^\ell$ with $\ell$ prime to the characteristic, and $X=\mathbf P^1_{R'}$ viewed over $R$. It is a regular projective flat surface with connected fibre $X_s=\ell\mathbf P^1_k$. The degree-zero Picard group of this thickened projective line is trivial (successive nilpotent quotients have $H^1=0$), hence unipotent, but its multiplicity gcd is $\ell$. The generic constants are $K(u)$, explaining exactly why the preceding order argument fails. \hfill$\square$

\exercise{4.6}{Compute pushforward functions and test exactness of the generic-fibre Picard sequence for a projective morphism with geometrically integral fibres.}
\pf (a) The proper geometrically integral generic fibre has constant field $k(\xi)$. The Stein algebra $f_*\mathcal O_X$ is finite over $\mathcal O_Y$ and has that same generic field. Normality of $Y$ therefore gives $f_*\mathcal O_X=\mathcal O_Y$. The projection formula implies that a line bundle whose pullback is trivial was already trivial, proving Picard injectivity.

\textbf{Part (b) requires flatness; the stated assumptions do not imply it.} Let $Y=\mathbf A^2_k$ and $X=\operatorname{Bl}_0Y$. Both schemes are regular, the map is dominant and projective, and every fibre is geometrically integral (a point or the exceptional $\mathbf P^1$). Its vertical prime exceptional divisor maps to a codimension-two point, not to a prime divisor. Moreover $\operatorname{Pic}(Y)=0$, $\operatorname{Pic}(X)=\Z[E]$ and $\operatorname{Pic}(X_\xi)=0$, so the proposed middle exactness fails.

With flatness added, the dimension formula makes the image of a vertical prime a codimension-one prime $E$ of $Y$. Geometric integrality and reducedness of the fibres make $f^{-1}E$ an integral Cartier divisor with multiplicity one; hence $D=f^*E$. A line bundle trivial on the generic fibre is represented by a vertical divisor and therefore descends from $Y$. Conversely every line bundle on the regular generic fibre is a Cartier divisor class; take its Weil-divisor closure in regular $X$ to extend it. These facts prove the corrected short exact Picard sequence. \hfill$\square$

\exercise{4.7}{Compute the component group by Smith normal form, evaluate it for $\mathrm I_n^*$, and prove blow-up invariance.}
\pf Retain the subsection's connected fibre and algebraically closed residue field, so $M$ has rank $n-1$ and primitive kernel $X_s/d$. The Euclidean algorithm on rows and columns gives unimodular $P,Q$ with
\[ PMQ=\operatorname{Diag}(a_1,\ldots,a_{n-1},0),\qquad
0<a_1\mid\cdots\mid a_{n-1}. \]
Row or column addition, exchange and sign change implement that algorithm. Since $\Phi(X)$ is the torsion subgroup of $\operatorname{coker}M$, this proves $\Phi(X)=\prod_{i<n}\Z/a_i\Z$.

For $\mathrm I_m^*$ the matrix is the negative affine $D_{m+4}$ Cartan matrix. The primitive null vector has a coefficient one at an end. Use it as one integral basis vector; the corresponding row and column vanish and the remaining block is the finite $D_{m+4}$ Cartan matrix. Its lattice is
\[ D_r=\{(x_1,\ldots,x_r)\in\Z^r:\sum x_i\text{ is even}\},\quad r=m+4. \]
Its discriminant group has four elements: for even $r$ the two spinor classes have order two, and for odd $r$ a spinor class has order four. Equivalently this follows by pairing with the integer difference vectors and the half-vector $(1/2,\ldots,1/2)$. Thus
\[ \Phi(\mathrm I_m^*)\simeq
\begin{cases}(\Z/2\Z)^2,&m\text{ even},\\
\Z/4\Z,&m\text{ odd}.
\end{cases} \]

For a blow-up put $\mu_i=\mu_x(\Gamma_i)$. Replace the basis $\widetilde\Gamma_i,E$ by $\pi^*\Gamma_i=\widetilde\Gamma_i+\mu_iE,E$; this is unimodular. The projection formula makes its matrix $\operatorname{Diag}(M,-[k(x):k])$. In particular
\[ \widetilde\Gamma_i\cdot\widetilde\Gamma_j+
\mu_i(E\cdot\widetilde\Gamma_j)=\Gamma_i\cdot\Gamma_j. \]
Here $k$ algebraically closed gives $[k(x):k]=1$, and a row sign change makes the last block $1$. Cokernel torsion is therefore unchanged. Common regular refinements prove invariance for every regular model. If the algebraically closed convention is dropped, the block is $-[k(x):k]$, not $1$: blowing up a degree-$c$ point of $\mathbf P^1_R$ gives the matrix $\left(\begin{smallmatrix}-c&c\\c&-c\end{smallmatrix}\right)$ and torsion $\Z/c\Z$. Thus the unrestricted numerical claim would be false over a general residue field. \hfill$\square$

\exercise{4.8}{Show divisibility of necessary base-change ramification by multiplicities of components that survive in stable reduction.}
\pf (a) At $\xi$ the DVR valuation gives $\nu_\xi(t)=d$. At a dominating generic point $\xi'$ of the stable model, the reduced special fibre gives $\nu_{\xi'}(\tau)=1$, hence $\nu_{\xi'}(t)=e_{L/K}$. Restricting the valuation through the two DVRs gives
\[ e_{L/K}=e_{\mathcal O_{\mathcal C',\xi'}/\mathcal O_{\mathcal C,\xi}}\,d. \]
Thus $d\mid e_{L/K}$. The stray $\mathcal C'_s$ in the question's initial definition of $\Gamma$ should be $\mathcal C_s$.

(b) Remove from $\mathcal C$ the unwanted components and, if necessary, finitely many smooth boundary points. The reduced special open has projective completion the given stable union. Proposition 4.37 supplies a finite surjective map from a union of components of the potential stable fibre onto it. Each of its component generic points therefore has a dominating DVR as in (a). Applying (a) to every retained $\Gamma_i$ proves $d_i\mid e_{L/K}$.

(c) In a regular normal-crossings model $\Gamma$ is smooth. If $p_a(\Gamma)\ge2$, apply the stable version of Proposition 4.37 to its smooth projective completion. If its genus is one, use that proposition's explicit smooth-genus-one alternative. In both cases a component survives in the potential stable fibre and (a) applies. This is why an arbitrary rational component, which might disappear entirely, has no such divisibility guarantee. \hfill$\square$

\exercise{4.9}{Bound the maximal Abelian and toric ranks of normal models, prove the graph inequality, and show that the ranks cannot decrease under finite extension.}
\pf (a) Resolve any normal model. Exercise 3.17(a), for its connected exceptional fibres, shows that both ranks can only increase. All regular models have the same ranks: point blow-ups add only a rational exceptional tree, so both changes are zero, and a common refinement compares any two. Thus the maxima exist and equal the ranks on a regular model.

They are finite and in fact their sum is at most $g$. To see this directly, set $F=X_s/d$ on a regular model. The full primitive-fibre calculation of Exercise 9.1.14 injects $H^0(F,\mathcal O_F)$ into that of its connected reduced fibre, which over algebraically closed $k$ is $k$. Thus $h^0(F)=1$ and
\[ h^1(F)=1+(g-1)/d. \]
Restriction surjects onto $H^1((X_s)_{\mathrm{red}},\mathcal O)$, which bounds $a+t$. For $g\ge1$ the displayed number is at most $g$. For $g=0$, integrality of its positive Euler characteristic forces $d=1$ and $h^1(F)=0$.

(b) The Betti number is $\beta=|E|-|V|+1$, with loops counted twice in valence, so
\[ 2(\beta-1)=\sum_v(\deg v-2). \]
Deleting a terminal vertex and its edge changes neither $\beta$ nor connectedness of the remaining core. After pruning, all target valences are at least two. The complete-edge lifting hypothesis implies that every vertex above $v$ has valence at least $\deg v$. The subgraph above the core is connected (path lifting and the complete bipartite lifting between consecutive fibres show this), and each target vertex has at least one preimage. Consequently
\[ \sum_{v_1}(\deg v_1-2)\ge
\sum_{v_2}|f^{-1}(v_2)|(\deg v_2-2)
\ge\sum_{v_2}(\deg v_2-2). \]
This proves the inequality. An even shorter proof of the question's stronger hypothesis is to select one lift of every vertex and an edge above each target edge joining the selected lifts. These form a subgraph isomorphic to $G_2$, so 1.18 applies. The valence proof also works under the weaker local lifting of incident edges, whenever each connected component over the core maps onto it.

(c) Normalize a regular model after the finite base extension and then resolve. Each original component has a preimage component, and the normalization of that component maps finitely onto its original normalization. Over algebraically closed $k$ the purely inseparable part is a Frobenius twist, which preserves genus; the separable part has genus at least the target's whenever that genus is positive, by Hurwitz. Summing shows that Abelian rank cannot decrease, and resolution can only add a nonnegative contribution.

For the toric rank use the branch-incidence graph, including its singular-point vertices before resolving. First contract edges mapping into the interior of a target component; this collapses connected subgraphs and cannot increase the upstream Betti number. Each resulting component vertex contains a component finite and surjective over the target component, so every incident target marked point has a lift there. At a vertex above a target singular point, going down for the finite normal extension gives a lift of every incident target branch through that chosen point. Thus over the pruned target core the resulting incidence map has the local edge-lifting property in the preceding valence proof; each connected component maps onto the core by path lifting. That proof gives $\beta_{\mathrm{up}}\ge\beta_{\mathrm{down}}$. Purely inseparable normalization is a universal homeomorphism and preserves the geometric branches, so causes no loss either. Contracting connected exceptional subgraphs cannot increase rank in the downward direction, by 3.17(a). Therefore
\[ a(C_L)\ge a(C),\qquad t(C_L)\ge t(C). \]
Only the local branch-lifting assertion is needed here; finite base change need not give every possible pair of component lifts a common intersection. This avoids an unjustified complete-bipartite assertion for an arbitrary covering. \hfill$\square$



\end{document}
